\section{Builder}
\label{sec:pat-builder}

\problemstatement{Separate the construction of a complex object from its
representation, so the same step-by-step process can build different
configurations -- especially useful when an object has many optional
parameters.}

\begin{patternbox}[When You Reach for It]
The triggers are \emph{``an object with many optional fields''} (the
telescoping-constructor smell: \texttt{Pizza(true,false,true,...)}) and
\emph{``build a complex object in steps.''} When a constructor has a long
list of parameters, half of them optional, a Builder makes construction
readable and safe.
\end{patternbox}

\subsection*{Understanding the pattern}

Instead of one giant constructor, a Builder exposes chainable setter
methods (\texttt{addCheese()}, \texttt{addOlives()}) that accumulate
configuration, then a \texttt{build()} that produces the finished,
immutable product. Each setter returns the builder itself so calls read
like a sentence. The product's own constructor can stay private, taking
only the builder, guaranteeing objects are never half-initialised.

\begin{ideabox}[Chainable Steps, One Final Build]
Name each optional step explicitly and let the terminal \texttt{build()}
assemble the product. This turns an unreadable positional-argument call
into self-documenting, order-independent construction, and lets you
validate the whole configuration in one place before creating the object.
\end{ideabox}

\subsection*{C++ implementation}

\begin{lstlisting}
#include <bits/stdc++.h>
using namespace std;

class Pizza {
public:
    string size;
    bool cheese = false, olives = false, mushrooms = false;
    void describe() const {
        cout << size << " pizza"
             << (cheese ? " +cheese" : "")
             << (olives ? " +olives" : "")
             << (mushrooms ? " +mushrooms" : "") << "\n";
    }
};

class PizzaBuilder {
    Pizza p;
public:
    PizzaBuilder& setSize(const string& s) { p.size = s; return *this; }
    PizzaBuilder& addCheese()    { p.cheese = true;    return *this; }
    PizzaBuilder& addOlives()    { p.olives = true;    return *this; }
    PizzaBuilder& addMushrooms() { p.mushrooms = true; return *this; }
    Pizza build() { return p; }               // finished product
};

int main() {
    Pizza pizza = PizzaBuilder()
                    .setSize("Large")
                    .addCheese()
                    .addMushrooms()
                    .build();                 // reads like a sentence
    pizza.describe();
    return 0;
}
\end{lstlisting}

\begin{complexitybox}[Trade-offs]
\textbf{Pros.} Readable, order-independent construction; no telescoping
constructors; can enforce validation and immutability at \texttt{build}
time. \textbf{Cons.} More code than a plain constructor; overkill for
objects with a couple of fields. Worth it precisely when the parameter
list is long and largely optional.
\end{complexitybox}

\begin{takeawaybox}
Builder trades a long, error-prone constructor for a fluent chain of named
steps ending in \texttt{build()}. Recognise it by the smell of many
optional parameters or boolean flags at a call site. It is one of the most
frequently \emph{used} patterns in real code, not just interviews.
\end{takeawaybox}
